\chapter{Funções utilitárias de avaliação}\label{ap:evalutils}

O módulo \texttt{eval\_utils.py} concentra as funções reutilizadas por todos os experimentos do Capítulo~\ref{cap:investigacao}. As funções principais são reproduzidas abaixo.

\begin{lstlisting}[caption={Intervalo de confiança \textit{bootstrap} para ROC-AUC.},label={lst:bootstrap}]
def bootstrap_auc_ci(y_true, y_score, n_boot=1000,
                     alpha=0.05, seed=42):
    rng = np.random.default_rng(seed)
    n = len(y_true)
    aucs = np.empty(n_boot)
    for i in range(n_boot):
        idx = rng.integers(0, n, size=n)
        if len(np.unique(y_true[idx])) < 2:
            aucs[i] = np.nan
            continue
        aucs[i] = roc_auc_score(y_true[idx], y_score[idx])
    aucs = aucs[~np.isnan(aucs)]
    auc_point = roc_auc_score(y_true, y_score)
    lower = float(np.quantile(aucs, alpha / 2))
    upper = float(np.quantile(aucs, 1 - alpha / 2))
    return float(auc_point), lower, upper
\end{lstlisting}

\begin{lstlisting}[caption={Split temporal walk-forward 70/15/15.},label={lst:walkforward}]
def walk_forward_split(df, train_frac=0.70,
                       val_frac=0.15):
    n = len(df)
    n_train = int(n * train_frac)
    n_val = int(n * val_frac)
    train = df.iloc[:n_train]
    val = df.iloc[n_train:n_train + n_val]
    test = df.iloc[n_train + n_val:]
    return train, val, test
\end{lstlisting}

\chapter{Arquitetura TCN}\label{ap:tcn}

A TCN (\textit{Temporal Convolutional Network}) utiliza blocos de convolução causal dilatada com conexões residuais. A configuração TCN~[32,32] contém duas camadas com 32 filtros cada, \textit{kernel size} 3 e dilatações $[1, 2]$.

\begin{lstlisting}[caption={Bloco convolucional causal da TCN.},label={lst:tcn}]
class CausalConv1dBlock(nn.Module):
    def __init__(self, in_ch, out_ch, kernel_size,
                 dilation, dropout=0.2):
        super().__init__()
        padding = (kernel_size - 1) * dilation
        self.conv1 = nn.Conv1d(in_ch, out_ch,
            kernel_size, padding=padding,
            dilation=dilation)
        self.conv2 = nn.Conv1d(out_ch, out_ch,
            kernel_size, padding=padding,
            dilation=dilation)
        self.chomp = padding
        self.relu = nn.ReLU()
        self.dropout = nn.Dropout(dropout)
        self.downsample = (nn.Conv1d(in_ch, out_ch, 1)
                          if in_ch != out_ch else None)

    def forward(self, x):
        out = self.dropout(self.relu(
            self.conv1(x)[:, :, :-self.chomp]))
        out = self.dropout(self.relu(
            self.conv2(out)[:, :, :-self.chomp]))
        res = (self.downsample(x)
               if self.downsample else x)
        return self.relu(out + res)
\end{lstlisting}

